\Article{363}
Die Hauptaufgabe besteht darin, diese Reste zu bestimmen. Zuerst suche man den Rest, den die größte Periode bei der Division durch das Produkt aller zugehörigen Primfunktionen ergibt. Ist dieses Produkt
\[
\equiv x^\lambda-Ax^{\lambda-1}+\text{usw.}
\]
so ist dieser Rest $\equiv A$. Die Form dieses Produkts lässt sich leicht finden, und aus Kapitel \textsc{vi} folgt, dass $A=0$ ist, wenn $\nu$ durch ein Quadrat teilbar ist; andernfalls ist $A$ entweder $=+1$ oder $=-1$, je nachdem, ob die Anzahl der Primfaktoren von $\nu$ gerade oder ungerade ist.

Nun zerlege man diese größte Periode in kleinere Perioden und stelle die Glieder einer beliebigen Periode durch $x^{kp^\pi u}$ dar, sodass $k$ innerhalb jeder Periode eine
\SourceLeaf{242}
feste Zahl ist, zwischen den Perioden aber wechselt, während $\pi$ und $u$ innerhalb jeder Periode veränderlich sind und die Werte, die sie in einer Periode annehmen, auch in den übrigen annehmen können. Man nehme vorläufig eine beliebige Primfunktion $P$ als Basis an, und die Perioden $\Sigma x^{p^\pi u}$, $\Sigma x^{k'p^\pi u}$ usw. mögen bei der Division durch sie jeweils die Reste $A,A'$ usw. ergeben. Dann ist $\Sigma x^{p^\pi u}-A$ durch das Produkt aller Funktionen $(P,\rho^u)$ teilbar; $\Sigma x^{k'p^\pi u}-A'$ durch das Produkt aller Funktionen $(P,\rho^{k'u})$; usw., usw. Man sieht leicht, dass die Größen $A,A'$ usw. Wurzeln einer bekannten Kongruenz sind. Sind nämlich die Perioden der Wurzeln der Gleichung $x^\nu=1$, die den vorigen Perioden entsprechen, Wurzeln der Gleichung $Q=0$, so sind $A,A'$ usw. Wurzeln der Kongruenz $Q\equiv0$. Denn es gilt
\[
\begin{aligned}
A+A'+\text{usw.}&\equiv\text{der Summe der Perioden},\\
AA+A'A'+\text{usw.}&\equiv\text{der Summe der Quadrate der Perioden}
\end{aligned}
\]
usw., usw. Die Rechnung ist derjenigen völlig ähnlich, die wir in Kapitel \textsc{vi} dargestellt haben, wenn man für $\rho$ nun $x$ setzt, da man auch hier für $x^\nu$ die Einheit setzen kann, wie dort für $\rho^\nu$.

Hat man die Wurzeln $A,A'$ usw. gefunden, so wähle man eine als Rest der Periode $\Sigma x^{p^\pi u}$ und ordne danach die Reste der übrigen wie in Kapitel \textsc{vi}. Auch hier bleibt diese Wahl frei, da die Funktion $P$ bisher völlig unbestimmt ist. Die folgende Rechnung ist derjenigen ganz analog, die wir in Kapitel \textsc{vi} behandelt haben; jede Einzelheit darzustellen, würde uns zu weit führen. Ist man schließlich zu $\Sigma x^{p^\pi}$ gelangt, so ist die Hauptaufgabe erledigt. Denn setzt man
\[
P\equiv x^m+ax^{m-1}+bx^{m-2}+\text{usw.}
\]
so ist $-a\equiv\Sigma x^{p^\pi}$. Ebenso erhält man den zweiten Koeffizienten der übrigen Funktionen $(P,\rho^k)$, woraus sich die restlichen Koeffizienten von $P$ bestimmen lassen. Es kann häufig vorkommen, dass man zu identischen Kongruenzen gelangt, aus denen sich anscheinend nichts ableiten lässt. Wie man dieser Schwierigkeit begegnen kann, wird weiter unten gezeigt.

\Article{364}
Ein Beispiel wird dies alles viel deutlicher machen. Die Funktion $x^{15}-1$ soll modulo $17$ in Faktoren zerlegt werden. Es ist $m=4$, und da das Produkt aller elementaren Funktionen
\[
\equiv\frac{x^{15}-1\pdot x-1}{x^3-1\pdot x^5-1}=x^8-x^7+x^5-x^4+x^3-x+1
\]

\SourceLeaf{243}
ist, gibt es somit nur zwei Primfaktoren vierten Grades, $P$ und $P'$. Nun verteile man $x,xx,x^4,x^7,x^8,x^{11},x^{13},x^{14}$ auf diese beiden Perioden:
\[
\begin{aligned}
\Sigma x^{17^\alpha}&\equiv x+xx+x^4+x^8,\\
\Sigma x^{7\pdot17^\alpha}&\equiv x^7+x^{11}+x^{13}+x^{14}
\end{aligned}
\]
Modulo einer der Funktionen $P,P'$ sei
\[
\Sigma x^{17^\alpha}\equiv A,\quad\Sigma x^{7\pdot17^\alpha}\equiv A'
\]
Dann gilt
\[
\begin{aligned}
A+A'&\equiv1\\
AA&\equiv\Sigma x^{2\pdot17^\alpha}+\Sigma x^{3\pdot17^\alpha}+\Sigma x^{5\pdot17^\alpha}+\Sigma x^{9\pdot17^\alpha}\\
A'A'&\equiv\Sigma x^{14\pdot17^\alpha}+\Sigma x^{6\pdot17^\alpha}+\Sigma x^{5\pdot17^\alpha}+\Sigma x^{3\pdot17^\alpha}
\end{aligned}
\]
also
\[
\begin{aligned}
AA+A'A'&\equiv\Sigma x^{17^\alpha}+\Sigma x^{7\pdot17^\alpha}+4\Sigma x^{3\pdot17^\alpha}+2\Sigma x^{5\pdot17^\alpha}\\
&\equiv1-4-4\equiv-7
\end{aligned}
\]
Somit sind $A$ und $A'$ die Wurzeln der Kongruenz
\[
xx-x+4\equiv0\;(\mathrm{mod.}\,17)
\]
und diese sind $6,12$. Daher teilt $P$ den Ausdruck
\[
x^8+x^4+xx+x-6
\]
und ist
\[
\equiv x^4-6x^3-2xx-12x+1
\]
Weiter ist $P'$ dann $\equiv(P,\rho^7)$ und somit
\[
\equiv x^4-12x^3-2xx-6x+1
\]
